Despite static analysis and data visualization being among the oldest concepts in computer science, disappointingly few graphic designers have been significantly interested in the former. While plenty of tools exist for creating control flow graphs or visualizing the stack, the majority of these favor a ``function over form'' paradigm that focuses on creating lightweight, readable output, assuming the user is more concerned with the information than its presentation. However, this assumption does not always hold: presenting in any context, such as at a conference or school, generally favors quickly digestible and visually appealing content over static, large blocks of text. In these cases, the presenter is forced to choose between using the unappealing generated visualization and the slog of creating something presentable for themselves.

LLVManim was developed to address this problem. By pairing LLVM's powerful static analysis infrastructure with the Manim animation library, the tool automatically generates professional-quality visualizations from LLVM Intermediate Representation (IR) and runtime traces. Through a simple command-line interface, users can produce animated stack walkthroughs, IR spotlight views with Static-Single-Assignment (SSA) value tracking, Control Flow Graph (CFG), which represents the execution paths between basic blocks in a program, traversal animations, and machine-readable scene graph exports---all without writing any Manim code or manually post-processing compiler output.

Over the course of a single semester, we designed and built a three-layer pipeline (Ingestion, Transformation, and Rendering) comprising 26~Python modules and validated by 302~automated tests at 84\% line coverage. The remainder of this paper describes the tools and techniques that informed our design~(Sections~\ref{sec:related_work}--\ref{sec:manim}), the system architecture~(Section~\ref{sec:llvmanim}), the results of our implementation~(Section~\ref{sec:results}), lessons learned~(Section~\ref{sec:discussion}), and conclusions~(Section~\ref{sec:conclusion}).
